% Part: first-order-logic
% Chapter: tableaux
% Section: propositional-rules

\documentclass[../../../include/open-logic-section]{subfiles}

\begin{document}

\iftag{FOL}
      {\olfileid{fol}{tab}{prl}}
      {\olfileid{pl}{tab}{prl}}

\olsection{ప్రవచన నియమాలు}

\subsection{$\lnot$కు నియమాలు}

\begin{defish}
\AxiomC{\sFmla{\True}{\lnot !A}}
\RightLabel{\TRule{\True}{\lnot}}
\UnaryInfC{\sFmla{\False}{!A}}
\DisplayProof
\hfill
\AxiomC{\sFmla{\False}{\lnot !A}}
\RightLabel{\TRule{\False}{\lnot}}
\UnaryInfC{\sFmla{\True}{!A}}
\DisplayProof
\end{defish}

\subsection{$\land$కు నియమాలు}

\begin{defish}\noindent
\AxiomC{\sFmla{\True}{!A \land !B}}
\RightLabel{\TRule{\True}{\land}}
\UnaryInfC{\sFmla{\True}{!A}}
\noLine
\UnaryInfC{\sFmla{\True}{!B}}
\DisplayProof
\hfill
\AxiomC{\sFmla{\False}{!A \land !B}}
\RightLabel{\TRule{\False}{\land}}
\UnaryInfC{$\sFmla{\False}{!A} \quad \mid \quad \sFmla{\False}{!B}$}
\DisplayProof
\end{defish}

\subsection{$\lor$కు నియమాలు}

\begin{defish}
\AxiomC{\sFmla{\True}{!A \lor !B}}
\RightLabel{\TRule{\True}{\lor}}
\UnaryInfC{$\sFmla{\True}{!A} \quad \mid \quad \sFmla{\True}{!B}$}
\DisplayProof
\hfill
\AxiomC{\sFmla{\False}{!A \lor !B}}
\RightLabel{\TRule{\False}{\lor}}
\UnaryInfC{\sFmla{\False}{!A}}
\noLine
\UnaryInfC{\sFmla{\False}{!B}}
\DisplayProof
\end{defish}

\subsection{$\lif$కు నియమాలు}

\begin{defish}
\AxiomC{\sFmla{\True}{!A \lif !B}}
\RightLabel{\TRule{\True}{\lif}}
\UnaryInfC{$\sFmla{\False}{!A} \quad \mid \quad \sFmla{\True}{!B}$}
\DisplayProof
\hfill
\AxiomC{\sFmla{\False}{!A \lif !B}}
\RightLabel{\TRule{\False}{\lif}}
\UnaryInfC{\sFmla{\True}{!A}}
\noLine
\UnaryInfC{\sFmla{\False}{!B}}
\DisplayProof
\end{defish}

\subsection{కట్ నియమం}

\begin{defish}
\AxiomC{}
\RightLabel{\Cut}
\UnaryInfC{$\sFmla{\True}{!A} \quad \mid \quad \sFmla{\False}{!A}$}
\DisplayProof
\end{defish}

\Cut{} నియమాన్ని ముందున్న ఒక !!{signed formula} “కు” వర్తింపజేయరు;
బదులుగా, !!a{tableau}లోని ప్రతి శాఖను రెండుగా చీల్చి, ఒక శాఖలో
$\sFmla{\True}{!A}$ను, మరొక శాఖలో~$\sFmla{\False}{!A}$ను ఉంచనిస్తుంది.
ఇది అవసరం కాదు---సంవృత !!{tableau} ఉన్న ప్రతి !!{signed formula}s
సమితికి \Cut{} వాడని ఒక సంవృత టాబ్లో కూడా ఉంటుంది---కానీ
!!{tableau}sను అనుకూలంగా కలపడానికి ఇది వీలు కల్పిస్తుంది.

\end{document}
